% !TeX program = pdflatex
% Polityka bezpieczeństwa informacji --- Basilisk Systems P.S.A. (wzór uchwały Rady Dyrektorów).
% Polityka, o której mowa w § 14 ust. 3 umowy spółki, umowie akcjonariuszy i Umowie nr 2.
\documentclass[11pt,a4paper]{report}
\usepackage[T1]{fontenc}
\usepackage[utf8]{inputenc}
\usepackage[english,polish]{babel}
\usepackage{basilisk_i18n}
\usepackage{lmodern}
\usepackage{microtype}
\usepackage{setspace}
\usepackage{parskip}
\usepackage{needspace}
\usepackage{tabularx}

\BusinessPlanCompany{Basilisk Systems}
\BusinessPlanWatermarkText{Basilisk Systems}
\BusinessPlanWatermarkColor{black!6}
\BusinessPlanWatermarkFontSize{38}{46}
\BusinessPlanWatermarkAngle{38}
\BusinessPlanAuthorsText{Basilisk Systems --- Rada Dyrektorów}
\BusinessPlanConfidentialText{Dokument wewnętrzny}
\BusinessPlanFooterText{Basilisk Systems --- polityka bezpieczeństwa informacji}
\BusinessPlanVersionText{Wersja 1.0 --- projekt}

\setstretch{1.08}
\setlist[enumerate]{topsep=0.2em,itemsep=0.12em,parsep=0.05em,leftmargin=2.0em}
\setlength{\parindent}{0pt}
\setlength{\parskip}{0.45em}
\setlength{\emergencystretch}{4em}
\renewcommand{\arraystretch}{1.2}

\newcounter{paragraf}
\newcommand{\paragraf}[1]{%
  \refstepcounter{paragraf}%
  \Needspace{5\baselineskip}%
  \subsection*{\S~\theparagraf\ #1}%
}
\newenvironment{ContractClause}[1]{\paragraf{#1}}{\par}

\newcommand{\field}[1]{\textbf{[#1]}}
\newcommand{\CompanyFull}{Basilisk Systems prosta spółka akcyjna}

\hypersetup{
  pdftitle={Basilisk Systems P.S.A. --- polityka bezpieczeństwa informacji},
  pdfauthor={Basilisk Systems --- Rada Dyrektorów}
}

\begin{document}

\BusinessPlanTitlePage
  {Basilisk Systems}
  {Wersja 1.0 --- projekt}
  {Polityka bezpieczeństwa informacji}
  {Polityka bezpieczeństwa informacji Basilisk Systems P.S.A.}
  {Klasyfikacja, dostęp, udostępnianie, incydenty i procedura przy odejściu}
  {1.0}

\begin{center}
\textbf{Załącznik do uchwały nr \field{numer} Rady Dyrektorów \CompanyFull{} z dnia \field{data}}
\end{center}

\begin{ContractClause}{Cel i zakres}
\begin{enumerate}
  \item Polityka określa, jak Spółka chroni swoje informacje. Jest polityką bezpieczeństwa informacji, o której mowa w §~14 ust.~3 umowy spółki, w umowie akcjonariuszy i w umowach przeniesienia praw, oraz niezbędnym działaniem w celu utrzymania poufności w rozumieniu art.~11 ust.~2 ustawy o zwalczaniu nieuczciwej konkurencji.
  \item Polityka wiąże dyrektorów, akcjonariuszy mających dostęp do informacji Spółki, pracowników i współpracowników (\textbf{„Użytkownicy”}). Każdy Użytkownik potwierdza zapoznanie się z nią przed uzyskaniem dostępu.
  \item Za wykonanie polityki odpowiada dyrektor wyznaczony przez Radę Dyrektorów (\textbf{„Odpowiedzialny”}), a do czasu wyznaczenia --- każdy dyrektor.
\end{enumerate}
\end{ContractClause}

\begin{ContractClause}{Klasyfikacja i oznaczanie}
\begin{enumerate}
  \item Informacje dzieli się na: \textbf{Publiczne} (opublikowane przez Spółkę); \textbf{Wewnętrzne} (organizacja pracy, informacje ogólne); \textbf{Poufne} (kod, modele, dane, dokumentacja techniczna, wyniki testów, informacje handlowe i finansowe, księga akcyjna, warunki finansowania, dane klientów i partnerów); \textbf{Kontrolowane} (Poufne, które dotyczą technologii uzbrojenia albo podwójnego zastosowania, albo pochodzą od klientów z sektora obronnego).
  \item Informacja nieoznaczona, niebędąca Publiczną, jest Poufna. Dokumenty Poufne i Kontrolowane oznacza się napisem „POUFNE” albo „KONTROLOWANE”, a repozytoria i katalogi --- w opisie; brak oznaczenia nie znosi poufności.
\end{enumerate}
\end{ContractClause}

\begin{ContractClause}{Kto ma dostęp do czego}
\begin{enumerate}
  \item Dostęp przyznaje się według zasady niezbędnej wiedzy, w zakresie zadań Użytkownika:

\begin{center}
\begin{tabularx}{\linewidth}{@{}>{\raggedright\arraybackslash}p{2.6cm} >{\raggedright\arraybackslash}X >{\raggedright\arraybackslash}X@{}}
\toprule
\textbf{Kategoria} & \textbf{Kto ma dostęp} & \textbf{Warunki} \\
\midrule
Publiczne & każdy & --- \\
Wewnętrzne & Użytkownicy & potwierdzenie zapoznania się z polityką \\
Poufne & Użytkownicy w zakresie zadań; osoby trzecie w zakresie Celu NDA & zgoda Odpowiedzialnego; dla osób trzecich --- uprzednio zawarta NDA według wzoru Spółki albo ustawowa tajemnica zawodowa \\
Kontrolowane & wyłącznie osoby z imiennej listy zatwierdzonej przez Radę Dyrektorów & tylko na urządzeniach Spółki; przekazanie za granicę albo osobie z państwa spoza UE --- po ocenie wymogów kontroli eksportu i uzyskaniu wymaganego zezwolenia \\
\bottomrule
\end{tabularx}
\end{center}

  \item Odpowiedzialny prowadzi rejestr dostępów (osoba, systemy i repozytoria, kategoria, data nadania i odebrania) i co kwartał sprawdza, czy dostępy odpowiadają zadaniom.
  \item Spółka nie przetwarza informacji niejawnych w rozumieniu ustawy o ochronie informacji niejawnych. Przed zawarciem umowy wymagającej dostępu do takich informacji (w szczególności z Ministerstwem Obrony Narodowej albo podmiotem mu podległym) Rada Dyrektorów wszczyna postępowanie o wydanie świadectwa bezpieczeństwa przemysłowego i wyznacza pełnomocnika do spraw ochrony informacji niejawnych.
\end{enumerate}
\end{ContractClause}

\begin{ContractClause}{Zasady techniczne}
\begin{enumerate}
  \item Informacje Wewnętrzne, Poufne i Kontrolowane przetwarza się wyłącznie na kontach Spółki, z uwierzytelnianiem dwuskładnikowym i hasłami przechowywanymi w menedżerze haseł. Kod przechowuje się w prywatnych repozytoriach w organizacji Spółki; zakazane jest kopiowanie go do prywatnych repozytoriów i kont.
  \item Urządzenia używane do pracy mają szyfrowanie dysku, blokadę ekranu i aktualne oprogramowanie. Informacje Kontrolowane przetwarza się tylko na urządzeniach Spółki.
  \item Informacji Poufnych i Kontrolowanych nie wprowadza się do ogólnodostępnych narzędzi i usług, w tym narzędzi sztucznej inteligencji, chyba że Odpowiedzialny zatwierdził narzędzie zapewniające, że dane nie są przechowywane ani wykorzystywane przez dostawcę.
  \item Klucze, tokeny i sekrety przechowuje się w menedżerze sekretów, a nie w kodzie ani w wiadomościach.
\end{enumerate}
\end{ContractClause}

\begin{ContractClause}{Udostępnianie na zewnątrz}
Informacje Poufne udostępnia się osobom trzecim (inwestorom, partnerom, doradcom, podwykonawcom) wyłącznie po zawarciu NDA według wzoru Spółki albo umowy z równoważną klauzulą, w zakresie niezbędnym do jej celu. Odpowiedzialny prowadzi ewidencję zawartych NDA i przekazanych materiałów. Informacje Kontrolowane udostępnia się osobom trzecim wyłącznie za zgodą Rady Dyrektorów.
\end{ContractClause}

\begin{ContractClause}{Incydenty}
Każdy Użytkownik niezwłocznie, nie później niż w ciągu 24 godzin, zgłasza Odpowiedzialnemu utratę urządzenia, podejrzenie nieuprawnionego dostępu albo ujawnienia informacji. Odpowiedzialny dokumentuje incydent, odbiera albo zmienia zagrożone dostępy i sekrety, a w razie naruszenia ochrony danych osobowych zapewnia zgłoszenie organowi nadzorczemu w terminie 72 godzin (art.~33 RODO).
\end{ContractClause}

\begin{ContractClause}{Procedura przy odejściu}
\begin{enumerate}
  \item Przy ustaniu współpracy z Użytkownikiem, z jakiejkolwiek przyczyny:
  \begin{enumerate}[label=\alph*)]
    \item najpóźniej w ostatnim dniu współpracy, a w razie sporu --- niezwłocznie, Odpowiedzialny odbiera wszystkie dostępy (konta, repozytoria, chmura, VPN, klucze i tokeny) i zmienia sekrety współdzielone, które Użytkownik znał;
    \item Użytkownik zwraca urządzenia, nośniki i dokumenty Spółki; zwrot potwierdza protokół;
    \item w terminie 14 dni Użytkownik wydaje wszystkie utwory i składa końcowe potwierdzenie, jeżeli wymaga tego jego umowa, oraz usuwa informacje Spółki z urządzeń i kont prywatnych;
    \item Użytkownik podpisuje oświadczenie według wzoru poniżej; Odpowiedzialny przypomina mu na piśmie o obowiązkach poufności oraz --- jeżeli dotyczą --- o zakazie konkurencji i zakazie nakłaniania;
    \item Odpowiedzialny aktualizuje rejestr dostępów, a w razie podejrzeń przegląda dzienniki dostępu z ostatnich 30 dni.
  \end{enumerate}
  \item Odmowa złożenia oświadczenia nie zwalnia z obowiązków; Odpowiedzialny odnotowuje ją w rejestrze.
\end{enumerate}
\end{ContractClause}

\begin{ContractClause}{Postanowienia końcowe}
Naruszenie polityki może stanowić podstawę rozwiązania umowy z Użytkownikiem oraz roszczeń, w tym o zapłatę kar umownych przewidzianych w umowach. Rada Dyrektorów przegląda politykę co najmniej raz w roku. Polityka wchodzi w życie z dniem \field{data}.
\end{ContractClause}

\Needspace{14\baselineskip}
\subsection*{Wzór oświadczenia przy odejściu}

Ja, \field{imię, nazwisko}, w związku z zakończeniem współpracy z \CompanyFull{} z dniem \field{data} oświadczam, że: (1)~zwróciłem/zwróciłam wszystkie urządzenia, nośniki i dokumenty Spółki; (2)~wydałem/wydałam Spółce wszystkie utwory i materiały powstałe w ramach współpracy; (3)~usunąłem/usunęłam informacje Spółki z urządzeń i kont prywatnych, z wyjątkiem \field{brak / wykaz kopii wymaganych przepisami prawa}; (4)~nie zachowałem/zachowałam żadnych kopii kodu, modeli, danych ani dokumentacji Spółki. Przyjmuję do wiadomości, że obowiązki poufności oraz --- jeżeli dotyczą --- zakaz konkurencji i zakaz nakłaniania obowiązują mnie nadal w zakresie określonym w moich umowach ze Spółką i w przepisach prawa.

\field{miejscowość, data} \hfill \field{podpis}

\end{document}
